\documentclass[12pt, a4paper]{article}

% --- Paquetes de Idioma y Codificación ---
\usepackage[utf8]{inputenc}
\usepackage[T1]{fontenc}
\usepackage[spanish, es-tabla]{babel}

% --- Tipografía (Times New Roman) ---
% El paquete 'newtxtext' y 'newtxmath' proporciona una versión moderna de Times
\usepackage{newtxtext,newtxmath}

% --- Configuración de Color de Fondo ---
\usepackage{xcolor}
\definecolor{fondocrema}{RGB}{255, 253, 235} % Amarillo muy ligero
\pagecolor{fondocrema}

% --- Márgenes y Geometría ---
\usepackage[margin=2.5cm]{geometry}

% --- Enlaces e Índice ---
\usepackage{hyperref}
\hypersetup{
    colorlinks=true,
    linkcolor=blue,
    filecolor=magenta,      
    urlcolor=cyan,
    pdftitle={Plantilla con Fondo Amarillento},
}

% --- Título y Autor ---
\title{Plantilla Académica Estilo Clásico \\ \large Una Guía sobre el Formato de Artículo}
\author{Nombre del Autor}
\date{\today}

\begin{document}

\maketitle

\begin{abstract}
Este documento sirve como ejemplo de una plantilla en LaTeX con fondo ligeramente amarillento, diseñada para reducir la fatiga visual y dar un aspecto elegante. Incluye una estructura de cinco secciones y un índice automático.
\end{abstract}

\newpage
\tableofcontents
\newpage

\section{Introducción}
Esta es la primera sección de tu documento. LaTeX permite organizar el contenido de manera jerárquica y profesional. En esta plantilla hemos configurado la fuente \textit{Times New Roman} (mediante el paquete \texttt{newtxtext}) para mantener un estándar académico.

\section{Estado del Arte}
Aquí puedes describir las investigaciones previas. El color de fondo se ha definido mediante el paquete \texttt{xcolor} usando una mezcla RGB que imita el papel de alta calidad.

\section{Metodología Utilizada}
En esta sección se detallan los procedimientos. Gracias al uso del tamaño de papel A4 y márgenes de 2.5cm, el texto mantiene una legibilidad óptima. 
\begin{itemize}
    \item Recolección de datos.
    \item Procesamiento estadístico.
    \item Validación de resultados.
\end{itemize}

\section{Análisis de Resultados}
Presentación de los datos obtenidos. Puedes incluir tablas, figuras y ecuaciones matemáticas complejas que se renderizarán perfectamente con la tipografía Times:
\begin{equation}
    E = mc^2
\end{equation}

\section{Conclusiones y Recomendaciones}
Finalmente, se resumen los hallazgos. Esta plantilla es ideal para trabajos que requieren una lectura prolongada, ya que el tono amarillento es más suave para la vista que el blanco puro.

\end{document}